% Source: https://artofproblemsolving.com/community/c3h1389273_analytic_geometry
% Stable identity: aops-topic-1389273-q01
The first graph is the union of the lines $\ell_1: y=x+2$ and $\ell_2: y=4-3x$; the second is the union of $\ell_3: y=2-x$ and $\ell_4: 2x-5y+7=0$. A common point lies on one line from each pair, so we intersect each of $\ell_1,\ell_2$ with each of $\ell_3,\ell_4$:
\begin{itemize}
\item $\ell_1\cap\ell_3$: $x+2=2-x$, so $(0,2)$;
\item $\ell_1\cap\ell_4$: $2x-5(x+2)+7=0$, so $(-1,1)$;
\item $\ell_2\cap\ell_3$: $4-3x=2-x$, so $(1,1)$;
\item $\ell_2\cap\ell_4$: $2x-5(4-3x)+7=0$, so $\left(\dfrac{13}{17},\dfrac{29}{17}\right)$.
\end{itemize}
No two of these lines are parallel, and the four points are distinct.

\begin{asy}
usepackage("amsmath");
size(6.5cm);
defaultpen(fontsize(10pt));
pen p1=blue, p2=red;
draw((-2.2,-0.2)--(1.2,3.2),p1);
draw((0.4,2.8)--(1.6,-0.8),p1);
draw((-1.2,3.2)--(2.2,-0.2),p2);
draw((-2.2,0.52)--(2.2,2.28),p2);
pair P1=(0,2),P2=(-1,1),P3=(1,1),P4=(13/17,29/17);
dot(P1);dot(P2);dot(P3);dot(P4);
label("$(0,2)$",P1,W,fontsize(8pt)); label("$(-1,1)$",P2,NW,fontsize(8pt)); label("$(1,1)$",P3,2*dir(15),fontsize(8pt));
label("$\left(\dfrac{13}{17},\dfrac{29}{17}\right)$",P4+(0.44,0.6),fontsize(8pt));
draw((-2.3,0)--(2.4,0),Arrow); draw((0,-0.9)--(0,3.4),Arrow);
label("\mbox{}",(min(currentpicture,user=true)+max(currentpicture,user=true))/2,black); // keeps the text after the figure black when LaTeX typesets the labels inline
\end{asy}

So the graphs have $\boxed{4}$ common points; the poster's answer is correct.
